\subsection{The level at the assigned prime and its dimension}\label{subsec:selected-word-level}
\hypertarget{source-selected-word-level}{}

Fix the two common split primes \(g_0,g_1\) and the weights of a point. The assigned prime varies with the word. Its contribution to the dimension must be explicit before the columns can be placed in a common rank range. For each sequence \(u\), choose the third prime \(T_u\) and the level there as in Proposition~\ref{prop:assigned-prime}. Write \(F=E^+\), \(f=[F:\mathbb Q]\), and use \(M_j\)
for the self-dual quotient volume \eqref{eq:row-self-dual-mass} in
rank \(j\). We use the class and volume calculation of the companion
reconstruction \cite{companion-word-rank}.

\subsubsection{The class count in the marked genus}

The marking changes the local stabilizer at the assigned prime. Counting its nonzero possibilities in the split genus isolates the entire dependence of the column dimension on that prime.

\begin{proposition}[Dimension of the marked definite column]\label{prop:marked-column-dimension}
For the column lattice of discriminant index \(T_u^{2f}\) at the
split assigned prime, with the nonzero marking at each place over
\(T_u\), the definite column has dimension

\[
 d_u=d_{u,*}(T_u^k-1)^f.
\]
 The factor \(d_{u,*}\) is computed
from the self-dual class volume, the common index and the
definite weights, before the marking is imposed.
\end{proposition}

\begin{proof}
\proofstep{The self-dual enlargement.}

At \(v\mid T_u\) write the two split parts of the column lattice as
\(B_+\oplus B_-\). Their duals for the cross pairing are
\(B_-^*\oplus B_+^*\). The marking is a nonzero residue row on the
discriminant module with scalars through one residue factor, called
the plus factor.

At \(v\mid T_u\) write \(l_v^+,l_v^-\) for the two defect lengths.
Elementary divisors for the cross pairing identify them. The nonzero
marking and the discriminant index \eqref{eq:selected-column-discriminant}
give
\[
 l_v^+=l_v^-\ge1,\qquad
 \sum_{v\mid T_u}(l_v^++l_v^-)=2f .
\]
There are \(f\) places, so both lengths are one at each, and \(B_+\) is the kernel of the row
on \(B_-^*\). Take \(N_v=B_-^*\oplus B_-\) at these places and the
original self-dual lattice elsewhere. Taking the inverse image of
\(\bigoplus_v N_v/B_v\) in \(B^\vee\) gives a global lattice \(N\).
Duality commutes with localization, hence \(N=N^\vee\).

Write \(K\) for its finite stabilizer. An isometry preserving the
marked lattice preserves \(N\) by its construction. In \(K_v\) at
the prime assigned to the sequence, fixing the residue row preserves its kernel and
hence the original lattice. The level index is \(T_u^k-1\) there
by transitivity on nonzero rows and surjectivity of reduction. At
each \(g_i\) the fine level reduces to the upper unitriangular group
in \(\mathrm{GL}_k\). Put
\[
 I_k=\prod_{i=0,1}
 \left(\frac{|\mathrm{GL}_k(\mathbf F_{g_i})|}
 {g_i^{k(k-1)/2}}\right)^f .
\]
At other places the index is one.

For a coarse double coset \(y\),
the fine classes above it are the orbits of
\(\Delta_y=y^{-1}U(F)y\cap K\) on the finite coset space of the
fine level in \(K\). Each orbit has \(|\Delta_y|\) elements by the
pro-\(g_i\) and cyclotomic degree argument of
Section~\ref{sec:ramified-column}. The coarse reciprocal-stabilizer
sum is \(M_k\), since \(E/F\) is unramified at the finite places
and \(N\) is self-dual. Summing the orbit counts gives
\[
 \#\{\text{fine classes}\}
 = I_k(T_u^k-1)^f\sum_y |\Delta_y|^{-1}
 = M_k I_k(T_u^k-1)^f .
\]

For each such class the selected column reads a form at its
representative in the frame coordinates of
Section~\ref{subsec:selected-column-metric}. The class
parametrization is bijective and changing the real frames is
invertible. Trivial stabilizers imply that any given collection of
coefficient values extends uniquely by equivariance. The local
coefficient space is the polynomial Weyl space of the recorded
weight with its determinant twist. Evaluation in a frame preserves
its dimension since frames form a dense open set after
complexification. If \(a_{\rho,i}\) are its shifted highest-weight
coordinates, put
\[
 V_{u,\rho}=\prod_{i<j}(a_{\rho,i}-a_{\rho,j}),\qquad
 F_j=\prod_{i=0}^{j-1}i!,\qquad
 d_{u,*}=M_k I_k\prod_\rho(V_{u,\rho}/F_k).
\]
Weyl's formula and the free choice of values give
\(d_u=d_{u,*}(T_u^k-1)^f\). The same formula uses the weight of the
finite class lattice, since adding a constant to every coordinate
leaves the differences unchanged.
\label{cor:selected-word-level-1}

\end{proof}

\subsubsection{The target before prime selection}

After removing the local marking factor from the column dimension, the remaining data specify a real target for the prime. Its size and variation across words determine how accurately a prime must approximate it.

Write \(r=k+m-1\), \(g=g_0g_1\), \(L=\log b=x^\epsilon\),
\(v=d_E/2\) with \(d_E\) the root discriminant, and let
\(X_u=(D_*/d_{u,*})^{1/(fk)}\).

\begin{proposition}[Bounds for the target before selection]\label{prop:prime-target-size}
Let \(d_{u,*}\) be the pre-level column factor of
Proposition~\ref{prop:marked-column-dimension}, and \(D_*\) the
identity term of Proposition~\ref{prop:projector-dimension}.
Set \(X_u=(D_*/d_{u,*})^{1/(fk)}\), \(g=g_0g_1\), \(v=d_E/2\).
Then
\[
 \log X_u\ge \sum_j f^{-1}\sum_\rho\log(w_{j\rho}v/\pi)+m\log g-O(L),
 \qquad \log X_u\ll(1+C_g)mL.
\]
Consequently \(t=X_u\Gamma^{-2/(fk)}\) satisfies \(\log t\ge mL/2\)
for every word at one sufficiently late threshold, using
\(w_{j\rho}\ge1\), \(g_i\ge(kb)^8\),
\(m\le x^{1/2+\epsilon/100}\) and \(k\ge m^5 Lx^2\).
\end{proposition}

\begin{proof}
\proofstep{Compare the level indices and formal degrees.}

Both quantities in this ratio
are defined before \(T_u\) is chosen. Let \(I_r\) be the original row-level index and \(d_\rho\) the
formal degree in compact-transfer measure. The identity term of the
trace and Harish--Chandra's formal degree formula give
\[
 D_*=M_r I_r\prod_\rho d_\rho,\qquad
 d_\rho F_r=V_{\lambda,\rho}
\]
for the regular row parameter.
The compact scalar circle has mass one and its central character
has absolute value one in the squared matrix coefficient.
\cite{atiyah-schmid}

Put \(s=m/2-1\). In each row reduction over \(g_i\) there are \(s\)
single blocks on either side of the block of size \(k+1\).
All strictly upper off-block entries are free. The middle block
is \(\operatorname{diag}(A,1)\) with \(A\) upper unitriangular.
Its subgroup together with those entries has order
\(g_i^{(r^2-(k+1)^2-2s)/2+k(k-1)/2}\). Dividing its index by the
column index and using the finite general linear orders gives
\[
 I_r/I_k=g^{f(km+m(m-1)/2)}
       \prod_{i=0,1}\prod_{j=k+1}^r(1-g_i^{-j})^f .
\]

In \eqref{eq:row-self-dual-mass} the ratio \(M_r/M_k\) cancels
all Euler factors up to \(k\), and the factorials cancel those of
the degree and dimension formulas. Use \(L_j\) for the Euler
products of \eqref{eq:row-self-dual-mass}, and put
\(J=(m-1)k+m(m-1)/2\). This gives the identity
\begin{align*}
 fk\log X_u={}& fJ\log(v/\pi)-\tfrac12 f(m-1)\log d_E\\
 &+f(km+m(m-1)/2)\log g
       +\sum_\rho\log(d_\rho F_r/V_{u,\rho})+Z,\\
 Z={}&\sum_{j=k+1}^r\log L_j+
       f\sum_{i=0,1}\sum_{j=k+1}^r\log(1-g_i^{-j}).
\end{align*}
For \(j\ge2\), the Euler-product majorant gives
\(|\log L_j|\le f\log\zeta(j)\), and
\(-\log(1-g_i^{-j})\le\log\zeta(j)\). Thus
\[
 |Z|\le 3f\sum_{j=k+1}^r\log\zeta(j)\ll fm .
\]

\proofstep{Compare the uncancelled Vandermonde factors.}

After cancellation, the remaining weight differences control variation of the target across words.


At a real place cancel the Vandermonde on the \(k-2\) interior
nodes against the one on the corresponding row nodes. The
logarithm of their ratio is \(O(k(1+\log k))\), by the binary
displacements of Section~\ref{sec:local-amplitudes}. At most \(2k-3\)
uncancelled column gaps involve an end coordinate, each at most
\(2a\) on the weight window. The cross gaps in the row contribute
at least \(\prod_j B_j^{k-1}\). Unused row gaps are at least one.
Here \(B_j\) is the absolute value of the \(j\)-th negative
coordinate and \(w_{j\rho}\) the corresponding chain weight.
The window has \(0<w_{j\rho}-B_j\le B_j/x\) and \(w_{j\rho}\le2b\).
Consequently
\[
 (k-1)\log B_j\ge k\log w_{j\rho}-k/x-\log(2b),\qquad
 \log(d_\rho F_r/V_{u,\rho})\ge k\sum_j\log w_{j\rho}-O(kL).
\]

We used \(m\le x,\ m\le k,\ \log k+\log(2a)=O(L)\).
In the opposite direction at most \((m+1)r\) row gaps remain,
each \(O(b^2)\), and each remaining column gap is at least one.
This bounds the log ratio from above by \(O(mkL)\).

\proofstep{Bound the remaining field and scale factors.}

The host contains the real subfield of the sixteenth cyclotomic
field. The conductor-discriminant formula there gives root
discriminant \(2^{11/4}>2\pi\). In the full host \(d_E\) is at most
its conductor \(N_E\). Thus \(0\le\log(v/\pi)\le L\) eventually,
and \(\log g\le 2C_gL\). Division by \(fk\) gives
\[
 \log X_u\ge
 \sum_j f^{-1}\sum_\rho\log(w_{j\rho}v/\pi)+m\log g-O(L),
 \qquad \log X_u\ll (1+C_g)mL .
\]
The replacement of \(J/k\) by \(m\) in the lower bound costs
at most \(L\). The gamma correction is \(2C_\Gamma m^5 L/k\).

In the chosen family \(m\le x^{1/2+\epsilon/100}\) and the lower
rank bound gives \(k\ge m^5 L x^2\) eventually
\label{cor:selected-word-level-2}. Thus this correction is \(o(1)\).
As \(w_{j\rho}\ge1\) and \(g_i\ge(kb)^8\), the target
\(t=X_u\Gamma^{-2/(fk)}\) has \(\log t\ge mL/2\).

\end{proof}

\subsubsection{Prime selection and uniform comparison}

The target has been bounded uniformly in the word. We must now realize it by a prime satisfying the splitting and congruence conditions closely enough to preserve the required dimension interval.

\begin{proposition}[Choice in the linked rank interval]\label{prop:assigned-prime}
For \(t=X_u\Gamma^{-2/(fk)}\) there is an assigned prime
\(T_u\in[t,(1+1/(fk))t]\) in the required splitting and
congruence class. With one threshold for the selected words, 
\[
 mL/2\le\log T_u\ll mL,\qquad
 \max\{T_u/T_{u'},T_{u'}/T_u\}\le \exp(1/2+O((1+\log k)/k)).
\]

\end{proposition}

\begin{proof}
\proofstep{A degree-one prime in the short interval.}

Use the identity ray class in \(\mathbb Q(\sqrt{-d})\) for
\(Q=N_Eg^2\). The relative Chebotarev estimate used for \(g_0,g_1\)
applies at \(t\) and \(t(1+1/(fk))\). Indeed
\(\log(4dQ^2)=O(L)\), \(\log t\gg mL\), and \(m\ge x^{1/4}\).
Its relative error is \(o(1/(fk))\) since \(fk\) grows polynomially.
In the exceptional-zero case put \(W(t)=1-t^{\beta-1}\), and
otherwise put \(W(t)=1\). Let \(t_1=t(1+1/(fk))\). For an exceptional zero, differentiation gives
\[
 \frac{d}{du}\operatorname{Li}(u^\beta)
 =\frac{\beta u^{\beta-1}}{\log(u^\beta)}
 =\frac{u^{\beta-1}}{\log u},\qquad
 \frac{d}{du}\big(\operatorname{Li}(u)-\operatorname{Li}(u^\beta)\big)
 =\frac{W(u)}{\log u}.
\] In the ordinary case the derivative of
\(\operatorname{Li}(u)\) is \(1/\log u\). As \(W\) increases and
\(t_1\le2t\), the main term increment before division by the
ray class number is at least
\[
 \int_t^{t_1} W(u)\,\frac{du}{\log u}
       \ge \frac{tW(t)}{fk\log(2t)} .
\]

The relative error bound at the two endpoints is
\(o(tW(t)/(fk\log t))\) in the same normalization.
For the exceptional case, write \(a=1-\beta>0\). Since
\(t\ge2\) and \(2t\le t^2\),
\[
 W(2t)\le1-t^{-2a}=(1+t^{-a})W(t)\le2W(t).
\]
For the exceptional zero in the Chebotarev estimate \(\beta\ge1/2\).
At either endpoint \(u\) use \(\log v\ge\beta\log u\) throughout the integral and \(\beta\ge1/2\): 
\[
 \int_{u^\beta}^{u}\frac{dv}{\log v}\le\frac{u-u^\beta}{\beta\log u}\le\frac{2uW(u)}{\log u}.
\]
 For the ordinary
case \(\operatorname{Li}(u)\ll u/\log u\) gives the same bound
with \(W=1\).
Write \(h_Q\) for the ray class number. The quadratic ideal
count gives \(h_Q=\exp(O(L))\), and Stark's bound gives
\(W(t)\ge\exp(-O(L))\).

For sufficiently large \(t\), the
two endpoint errors together are at most one half of the
main term lower bound. Thus the interval contains at least
\[
 \frac{tW(t)}{2fk\,h_Q\log(2t)}=t^{1-o(1)}
\]
prime ideals of the required class, since
\(\log t\gg mL\) and \(m\ge x^{1/4}\).
The \(O(\sqrt{2t}+\log Q)\) ideals of degree two or at
ramified residue primes cannot exhaust them. Choose a remaining degree-one prime ideal in the identity ray
class and a ray generator \(\pi_u=1+Qz\) with integral \(z\)
(the ideal is integral). In this quadratic field
\[
 T_u=N(\pi_u)=1+Q\operatorname{Tr}(z)+Q^2N(z)\equiv1\pmod Q .
\]
The norm congruence gives complete splitting in the host and the
required disjointness from the common level. Ratios of two such
generators give the norm scalars carrying the levels at the primes assigned to the sequences.
\cite[Theorems 1.4 and 3.3]{tz-chebotarev}

To compare two words at the same point, put \(J=k-2\).
At each real place the interior coordinates in
Definition~\ref{hyp:critical-level-weights} are binary perturbations
of a single arithmetic progression of spacing \(4\Delta\). The
logarithm of the interior Vandermonde differs from that of this
progression by at most \(J/4\) by Lemma~\ref{lem:binary-vandermonde-run}.
Using the same progression for both words gives
\[
 \left|\log\frac{V^{\mathrm{int}}_{u,\rho}}
                         {V^{\mathrm{int}}_{u',\rho}}\right|\le J/2 .
\]


The two endpoints agree for both words at this place. The gap from
the upper endpoint to any interior coordinate is at least
\(4k\Delta\). A gap to the lower endpoint at distance \(h\ge1\)
is at least \(3h\Delta\), also when that endpoint is replaced by
the lower top-place coordinate. Each interior coordinate moves by at
most \(\Delta\). Since \(|\log x-\log y|\le|x-y|/\min(x,y)\), the first
endpoint costs at most \(J/(4k)\), the last at most \(H_J/3\), with
\(H_J=\sum_{h=1}^J h^{-1}\). The endpoint pair itself contributes zero.
Thus the exact formula for the pre-level column and the prime interval give
\[
 \begin{aligned}
 |\log X_u-\log X_{u'}|
 &\le \frac{J/2+J/(4k)+H_J/3}{k},\\
 |\log T_u-\log T_{u'}|
 &\le |\log X_u-\log X_{u'}|+\log(1+1/(fk)).
 \end{aligned}
\]


In the first line the sum over the \(f\) real places cancels the
factor \(f\) in the exponent \(1/(fk)\). In the second, the common
factor \(\Gamma^{-2/(fk)}\) cancels and both logarithmic rounding
errors belong to the same interval of the displayed length. In
particular the ratio in either orientation is at most \(2\) at
one sufficiently late threshold. Proposition~\ref{prop:two-row-amplitude-bounds}
and Lemma~\ref{lem:row-corank-one-gain} can use \(A=2\) and the fixed
constant \(C_a=\max(1,C_{\rm pt},12e\pi)\). Increasing the same
threshold gives \(C_a^3\le k\).
\end{proof}
